\documentclass[11pt]{article}
\usepackage[letterpaper,margin=0.88in]{geometry}
\usepackage[T1]{fontenc}
\usepackage{lmodern}
\usepackage{microtype}
\usepackage{amsmath,amssymb,amsthm,mathtools,bm}
\usepackage{booktabs,array}
\usepackage{xcolor}
\usepackage[colorlinks=true,linkcolor=blue!48!black,citecolor=blue!48!black,urlcolor=blue!62!black]{hyperref}
\usepackage{enumitem}

\newtheorem{theorem}{Theorem}[section]
\newtheorem{proposition}[theorem]{Proposition}
\newtheorem{corollary}[theorem]{Corollary}
\newtheorem{lemma}[theorem]{Lemma}
\theoremstyle{definition}
\newtheorem{definition}[theorem]{Definition}
\theoremstyle{remark}
\newtheorem{remark}[theorem]{Remark}

\newcommand{\Tr}{\operatorname{Tr}}
\newcommand{\rank}{\operatorname{rank}}
\newcommand{\diag}{\operatorname{diag}}
\newcommand{\spec}{\operatorname{spec}}
\newcommand{\1}{\mathbf 1}
\newcommand{\norm}[1]{\left\lVert#1\right\rVert}
\newcommand{\HS}{\mathrm{HS}}
\newcommand{\ka}{\kappa_d}
\newcommand{\Am}{\mathcal A_m}
\allowdisplaybreaks

\title{\textbf{One-Spike Inverse Self-Commutators and Exact\\Three-versus-Four-Kick Curvature Synthesis}\\[1.5mm]
\large Singular rigidity, sharp stability, and a universal switching gain}
\author{Lluis Eriksson}
\date{30 August 2026}
\hypersetup{
 pdftitle={One-Spike Inverse Self-Commutators and Exact Three-versus-Four-Kick Curvature Synthesis},
 pdfauthor={Lluis Eriksson},
 pdfsubject={Exact inverse self-commutator costs and finite-switch matrix curvature synthesis},
 pdfkeywords={self-commutator, inverse commutator, Horn inequalities, Hilbert-Schmidt norm, matrix polygon, curvature synthesis}
}

\begin{document}
\maketitle

\begin{abstract}
For a traceless Hermitian matrix $F\in M_d(\mathbb C)$, let
\[
 \kappa_d(F)=\min\{\norm H_{\HS}\norm K_{\HS}:H=H^*,\ K=K^*,\ -i[H,K]=F\}.
\]
We solve this inverse problem on the complete spectral cone having one
positive eigenvalue.  If the nonzero spectrum is
$(P,-b_1,\ldots,-b_n)$ with $b_1\ge\cdots\ge b_n>0$ and
$\sum_jb_j=P$, then
\[
 \kappa_d(F)=\sum_{j=1}^n j b_j.
\]
Every balanced optimum has the same nonzero squared singular values
$2\sum_{j=\ell}^n b_j$, $1\le\ell\le n$, and hence rank $n$, independently
of ambient zero padding.  The deficit from the fixed-mass maximum is one half
of the pairwise dispersion of the $b_j$ and is sharply equivalent to their
$\ell^1$ distance from uniformity.

We then solve two fixed-target matrix-polygon problems in every finite
dimension.  For a balanced $m$-kick loop $x_1+\cdots+x_m=0$, define
\[
 H_{\rm geo}=\frac i2\sum_{k>j}[x_k,x_j],\qquad
 S_2=\sum_j\norm{x_j}_{\HS},
\]
and let $\mathcal A_m(F)$ be the infimum of $S_2^2$ subject to
$H_{\rm geo}=F$.  For every nonzero traceless Hermitian target,
\[
 \mathcal A_3(F)=12\sqrt3\,\kappa_d(F),\qquad
 \mathcal A_4(F)=16\,\kappa_d(F).
\]
The three-kick identity consolidates the author's earlier low-dimensional
cases, while the four-kick identity extends the previously known sign-paired
equality face to every target.  Thus a fourth kick lowers the exact optimal
squared perimeter by the target- and dimension-independent factor
$4/(3\sqrt3)$, a reduction of $1-4/(3\sqrt3)\approx23.02\%$.  The four-kick lower
bound uses an exact Gram-determinant argument; a tempting stronger
parallelogram inequality is false and is not used.  Exact-arithmetic replays
and an independent implementation accompany the manuscript.
\end{abstract}

\section{Question, answer, and scope}

Finite-dimensional self-commutator existence is classical
\cite{albert1957,thompson1958,fan1980}.  It does not determine the least
Hilbert--Schmidt resource of Hermitian factors generating a prescribed
target.  We study that inverse metric problem and its consequence for
balanced matrix loops.  Throughout,
\[
 \norm X_{\HS}^2=\Tr(X^*X)
\]
is unnormalized, $F=F^*$ is traceless, and scalar parts of Hermitian factors
are removed because they do not affect commutators.

The surrounding literature contains forward Frobenius commutator
inequalities \cite{bottcher2005,bottcher2008,cheng2010}, unrestricted-factor
inverse problems \cite{weiss1980}, and the weighted-shift constructions used
in self-commutator existence theory \cite{fan1980,beltita2014}.  The present
optimization keeps both factors Hermitian.  Its one-spike answer is extracted
from a small universal family of Horn inequalities rather than from a
dimension-by-dimension enumeration.

The loop problem is motivated by second-order product formulas and geometric
terms in controlled or Zeno dynamics \cite{burgarth2013,chen2022}.  Polygonal
isoperimetry itself is classical \cite{indrei2015}; our fixed-target quantity
retains the complete matrix $F$, not merely one norm of it.  Building on the
author's earlier dimension-three, dimension-four, and dimension-five inverse
formulas and loop corollaries \cite{erikssonQutrit,erikssonFour,erikssonFive},
on earlier balanced-kick mechanisms \cite{erikssonOperational}, and on an
earlier sign-paired matrix-polygon equality face
\cite{erikssonIsoperimetry}, the scoped package is:
\begin{enumerate}[label=(\roman*),leftmargin=2.2em]
 \item an exact all-dimensional formula on the one-positive or one-negative
 spectral cone;
 \item rigidity of the singular spectrum and rank of every balanced optimum;
 \item a sharp $\ell^1$ stability theorem at the uniform one-spike ray;
 \item a dimension-free consolidation of the triangular law and an all-target
 four-kick law extending the sign-paired case, including a strict universal
 switching gain.
\end{enumerate}
We do not claim novelty for the earlier low-dimensional slices, weighted-shift
self-commutator constructions, Horn feasibility, group-commutator loops, or
general Gini/mean-deviation inequalities.  We do not claim a closed formula for $\kappa_d$ when both sign multiplicities
exceed one, a classification of all optimizing matrices, an
infinite-dimensional result, or a general exact formula for five or more
kicks.  Priority is stated conservatively: the Horn and weighted-shift
machinery is standard, and the accompanying search is not a proof of novelty.

\section{Balanced self-commutators}

\begin{definition}
For traceless Hermitian $F\in M_d(\mathbb C)$ set
\begin{equation}
 \kappa_d(F)=\inf\{\norm H_{\HS}\norm K_{\HS}:H=H^*,\ K=K^*,\ -i[H,K]=F\}.
 \label{eq:kappa}
\end{equation}
\end{definition}

\begin{lemma}[one-matrix reduction]
\label{lem:one-matrix}
For every traceless Hermitian $F$,
\begin{equation}
 \kappa_d(F)=\frac12\min\{\norm C_{\HS}^2:CC^*-C^*C=2F\}.
 \label{eq:one-matrix}
\end{equation}
The minimum is attained.  Moreover, an optimal Hermitian pair may be chosen
so that
\begin{equation}
 \langle H,K\rangle_{\HS}=0,
 \qquad \norm H_{\HS}=\norm K_{\HS}=\sqrt{\kappa_d(F)}.
 \label{eq:orthogonal-optimum}
\end{equation}
\end{lemma}

\begin{proof}
If $F=0$, both minima are zero, attained by $H=K=C=0$, and the final
assertion is immediate.  Hence assume $F\ne0$.
If $-i[H,K]=F$, reciprocal scaling of $H,K$ preserves the commutator and
their norm product, and can make their norms equal.  For $C=H+iK$,
\[
 CC^*-C^*C=-2i[H,K]=2F,
 \qquad
 \norm C_{\HS}^2=\norm H_{\HS}^2+\norm K_{\HS}^2.
\]
Conversely, the Hermitian real and imaginary parts $X,Y$ of any feasible
$C$ obey $-i[X,Y]=F$ and
$\norm X_{\HS}\norm Y_{\HS}\le\norm C_{\HS}^2/2$.  This proves
\eqref{eq:one-matrix}.  For general feasibility, order an eigenbasis with all
positive eigenvalues first, then zeros, then negative eigenvalues.  Every
partial sum $s_j$ of this ordering is nonnegative and the final sum is zero.
The finite weighted shift $Ce_{j+1}=\sqrt{2s_j}\,e_j$ then satisfies
$(CC^*-C^*C)/2=F$.  Finite-dimensional compactness gives attainment.

Take a minimizing Hermitian pair and balance its nonzero norms.  Replacing
$K$ by
\[
 K-\frac{\langle H,K\rangle_{\HS}}{\norm H_{\HS}^2}H
\]
preserves the commutator.  Unless the inner product vanishes, it strictly
reduces the product after rebalancing, contradicting optimality.  This gives
\eqref{eq:orthogonal-optimum}.
\end{proof}

The Horn description of spectra of sums of Hermitian matrices
\cite{horn1962,klyachko1998,knutson1999,fulton2000} applies to
$CC^*+(-C^*C)=2F$.  We only need a nested family for which the
Littlewood--Richardson coefficient is visibly one.

\section{The complete one-spike cone}

Assume that $F$ has one positive eigenvalue.  Suppressing arbitrary zero
padding, write its nonzero spectrum as
\begin{equation}
 (P,-b_1,\ldots,-b_n),
 \qquad b_1\ge\cdots\ge b_n>0,
 \qquad \sum_{j=1}^n b_j=P.
 \label{eq:one-spike-spectrum}
\end{equation}
The order in \eqref{eq:one-spike-spectrum} is adapted to the constructor; the
decreasing spectral order is $(P,0,\ldots,0,-b_n,\ldots,-b_1)$.

\begin{theorem}[exact cost and optimizer rigidity]
\label{thm:one-spike}
For \eqref{eq:one-spike-spectrum},
\begin{equation}
 \boxed{\quad \kappa_d(F)=\sum_{j=1}^n j b_j.\quad}
 \label{eq:one-spike-cost}
\end{equation}
If $C=H+iK$ comes from any balanced optimal pair, then the nonzero squared
singular values of $C$, in decreasing order, are
\begin{equation}
 p_\ell=2\sum_{j=\ell}^n b_j,
 \qquad 1\le\ell\le n.
 \label{eq:rigid-spectrum}
\end{equation}
Consequently every such $C$ has rank exactly $n$.  The same statements hold,
with the same singular data, when the negative spectral subspace has rank one.
\end{theorem}

\begin{proof}
Let $d=n+1+z$, where $z$ is the zero multiplicity, and let
$p_1\ge\cdots\ge p_d\ge0$ be the eigenvalues of $CC^*$ for a minimizer in
\eqref{eq:one-matrix}.  If $p_d=t>0$, write the polar decomposition
$C=US^{1/2}$ with $S=C^*C$; invertibility makes $U$ unitary.  Then
$C'=U(S-t\1)^{1/2}$ satisfies
\[
 C'C'^*=CC^*-t\1,\qquad C'^*C'=C^*C-t\1.
\]
It preserves the self-commutator and strictly reduces the norm.  Hence
$p_d=0$ at an optimum.  The decreasing spectra in Horn's sum problem are
\[
 \alpha=(p_1,\ldots,p_{d-1},0),\quad
 \beta=(0,-p_{d-1},\ldots,-p_1),\quad
 \gamma=2\spec(F).
\]

For $1\le\ell\le n$, take
\[
 I_\ell=\{1,\ldots,\ell\},\qquad
 J_\ell=K_\ell=\{1,d-\ell+2,\ldots,d\}.
\]
These are valid Horn triples: the partition associated with $I_\ell$ is
zero and therefore
$c_{0,\lambda(J_\ell)}^{\lambda(K_\ell)}=1$.  The corresponding inequality
telescopes to
\begin{equation}
 p_\ell\ge 2\left(P-\sum_{j=1}^{\ell-1}b_j\right)
 =2\sum_{j=\ell}^n b_j.
 \label{eq:horn-tails}
\end{equation}
It follows that
\[
 \frac12\sum_{r=1}^d p_r
 \ge\sum_{\ell=1}^n\sum_{j=\ell}^n b_j
 =\sum_{j=1}^n j b_j.
\]

For attainment, use an eigenbasis ordered as in
\eqref{eq:one-spike-spectrum}, put
$t_\ell=\sum_{j=\ell}^n b_j$, and define the weighted shift
\begin{equation}
 Ce_{\ell+1}=\sqrt{2t_\ell}\,e_\ell
 \quad(1\le\ell\le n),
 \qquad C=0\quad\hbox{on the zero-padded complement}.
 \label{eq:weighted-shift}
\end{equation}
Direct multiplication gives
$(CC^*-C^*C)/2=F$ and
$\norm C_{\HS}^2/2=\sum_\ell t_\ell=\sum_j j b_j$.
Thus the lower bound is exact.  At any optimum, equality of the sum forces
equality in every nonnegative inequality \eqref{eq:horn-tails}, while all
remaining $p_r$ vanish.  This proves \eqref{eq:rigid-spectrum} and the rank
claim.  Replacing $F$ by $-F$ replaces $C$ by $C^*$.
\end{proof}

\begin{remark}[what rigidity does not say]
The theorem fixes the singular spectrum and rank of every balanced optimum;
it does not assert that every optimum is entrywise a weighted shift.  Unitary
and phase freedoms, and possibly further equality geometry, remain.
\end{remark}

\section{Sharp stability at the uniform ray}

For fixed $P$ and $n$, the general rank-adaptive upper endpoint is
$((n+1)/2)P$.  On the one-spike cone its deficit has an elementary dispersion
form.  Comparisons between empirical Gini mean difference and mean absolute
deviation are known \cite{cerone2006}; the theorem below records the exact
constants and equality families in this ordered, nonnegative, fixed-mass
spectral normalization.

\begin{theorem}[Gini identity and sharp trace-distance stability]
\label{thm:stability}
Let $u=P/n$ and
\[
 D=\frac{n+1}{2}P-\kappa_d(F),
 \qquad L=\sum_{j=1}^n|b_j-u|.
\]
Then
\begin{equation}
 \boxed{\quad
 D=\frac12\sum_{1\le i<j\le n}(b_i-b_j).
 \quad}
 \label{eq:gini}
\end{equation}
For $n\ge2$,
\begin{equation}
 \boxed{\quad \frac n4L\le D\le\frac{n-1}{2}L.\quad}
 \label{eq:stability}
\end{equation}
Both constants are optimal.  In particular, the uniform opposite-sign block
$b_1=\cdots=b_n=P/n$ is the unique fixed-mass maximizer of $\kappa_d$.
\end{theorem}

\begin{proof}
Substitution of \eqref{eq:one-spike-cost} gives
\[
 D=\frac12\sum_{j=1}^n(n+1-2j)b_j,
\]
which is \eqref{eq:gini}.  Put $y_j=b_j-u$ and let $k$ be the last index with
$y_k\ge0$.  The total positive and negative deviation masses both equal
$L/2$.  Keeping only the cross pairs in the nonnegative ordered sum gives
\[
 2D\ge\sum_{i\le k<j}(y_i-y_j)=\frac n2L,
\]
which is the lower bound.  Equality holds for every two-level deviation
profile: the $y_i$ are constant on $i\le k$ and on $i>k$.

For the upper bound, the coefficients
$w_j=(n+1-2j)/2$ satisfy $|w_j|\le(n-1)/2$, and hence
\[
 D=\sum_jw_jy_j\le\frac{n-1}{2}\sum_j|y_j|.
\]
Equality is obtained when only $y_1$ and $y_n$ are nonzero and opposite.
Finally $D=0$ exactly when every pairwise difference vanishes.
\end{proof}

The exact formula also orders the entire cone without solving any further
optimization problem.  Recall that a decreasing vector $b$ majorizes a
decreasing vector $c$ of the same mass if
$\sum_{j=1}^k b_j\ge\sum_{j=1}^k c_j$ for every $k<n$.

\begin{corollary}[strict Schur concavity on the one-spike cone]
\label{cor:majorization}
Let $b,c\in(0,\infty)^n$ be decreasing and have the same sum $P$, and let
$F_b,F_c$ be one-spike targets with opposite-sign blocks $b,c$.  If $b$
majorizes $c$, then
\[
 \kappa_d(F_b)\le \kappa_d(F_c).
\]
Equality holds if and only if $b=c$.  Thus making the opposite-sign spectral
mass more unequal strictly lowers the inverse cost, while the uniform block
is its unique fixed-mass maximum.
\end{corollary}

\begin{proof}
Writing $B_k=\sum_{j=1}^k b_j$ and using Theorem~\ref{thm:one-spike},
\[
 \kappa_d(F_b)=\sum_{j=1}^n j b_j
 =nP-\sum_{k=1}^{n-1}B_k.
\]
The corresponding identity for $c$ and the defining prefix inequalities give
the result.  Equality of the costs forces equality of every prefix, hence
$b=c$.
\end{proof}

\section{Exact finite-switch curvature synthesis}

Let $x_1,\ldots,x_m$ be Hermitian matrices with $\sum_jx_j=0$ and define
\begin{equation}
 H_{\rm geo}(x)=\frac i2\sum_{k>j}[x_k,x_j],
 \qquad S_2(x)=\sum_{j=1}^m\norm{x_j}_{\HS}.
 \label{eq:loop}
\end{equation}
For a traceless Hermitian target $F$, set
\begin{equation}
 \mathcal A_m(F)=\inf\{S_2(x)^2:\sum_jx_j=0,\ H_{\rm geo}(x)=F\}.
 \label{eq:action}
\end{equation}
Scalar components can be deleted from every edge without changing the target
and while weakly decreasing $S_2$, so the definition is already
gauge-minimal.

\subsection{Three kicks}

\begin{theorem}[exact triangular action]
\label{thm:three}
For every finite-dimensional traceless Hermitian $F$,
\begin{equation}
 \boxed{\quad \mathcal A_3(F)=12\sqrt3\,\kappa_d(F).\quad}
 \label{eq:three}
\end{equation}
\end{theorem}

\begin{proof}
The zero target is attained by the zero loop, so assume $F\ne0$.  Balance gives $x_3=-x_1-x_2$ and
$F=-i[x_1,x_2]/2$.  In the real Hilbert space of Hermitian matrices, remove
from $x_2$ its projection onto $x_1$.  The commutator is unchanged, so
\[
 \kappa_d(F)\le\frac12
 \sqrt{\norm{x_1}_{\HS}^2\norm{x_2}_{\HS}^2
       -\langle x_1,x_2\rangle_{\HS}^2}=:\Delta.
\]
Here $\Delta$ is the Euclidean area of the triangle with edge vectors
$x_1,x_2,x_3$.  The sharp triangular isoperimetric inequality gives
$S_2^2\ge12\sqrt3\,\Delta$, proving the lower bound.

Choose $H,K$ as in \eqref{eq:orthogonal-optimum}, put
$\alpha=2/\sqrt[4]{3}$, and set
\[
 x_1=\alpha H,\qquad
 x_2=\alpha\left(-\frac12H+\frac{\sqrt3}{2}K\right),\qquad
 x_3=-x_1-x_2.
\]
The three norms equal $\alpha\sqrt{\kappa_d(F)}$ and direct calculation gives
$H_{\rm geo}=F$.  Thus
$S_2^2=9\alpha^2\kappa_d(F)=12\sqrt3\,\kappa_d(F)$.
\end{proof}

\subsection{Four kicks and the Gram correction}

\begin{theorem}[exact quadrilateral action]
\label{thm:four}
For every finite-dimensional traceless Hermitian $F$,
\begin{equation}
 \boxed{\quad \mathcal A_4(F)=16\,\kappa_d(F).\quad}
 \label{eq:four}
\end{equation}
\end{theorem}

\begin{proof}
The zero target is attained by the zero loop, so assume $F\ne0$.  For any balanced four-tuple, put
\[
 D_1=x_1+x_2,\qquad D_2=x_2+x_3.
\]
A direct expansion gives the diagonal identity
\begin{equation}
 H_{\rm geo}=-\frac i2[D_1,D_2].
 \label{eq:diagonal-identity}
\end{equation}
Write
\[
 a=\norm{D_1}_{\HS}^2,\qquad
 b=\norm{D_2}_{\HS}^2,\qquad
 c=\langle D_1,D_2\rangle_{\HS},\qquad
 q=\sqrt{ab-c^2}.
\]
Shearing $D_2$ by its component parallel to $D_1$ leaves the commutator
unchanged, so \eqref{eq:diagonal-identity} implies
\begin{equation}
 \kappa_d(F)\le\frac q2.
 \label{eq:gram-kappa}
\end{equation}
Balance also gives
$D_1+D_2=x_2-x_4$ and $D_1-D_2=x_1-x_3$.  Therefore
\begin{align}
 S_2&\ge\norm{D_1+D_2}_{\HS}+\norm{D_1-D_2}_{\HS},\nonumber\\
 S_2^2&\ge
 2(a+b)+2\sqrt{(a-b)^2+4(ab-c^2)}\nonumber\\
 &\ge 8q\ge16\kappa_d(F).
 \label{eq:gram-chain}
\end{align}
The penultimate step uses $a+b\ge2q$ and
$\sqrt{(a-b)^2+4q^2}\ge2q$.

For attainment, choose the orthogonal equal-norm optimum $H,K$ from
\eqref{eq:orthogonal-optimum}, set $D_1=\sqrt2H$, $D_2=\sqrt2K$, and take
\begin{equation}
 x_1=\frac{D_1-D_2}{2},\quad
 x_2=\frac{D_1+D_2}{2},\quad
 x_3=-x_1,\quad x_4=-x_2.
 \label{eq:square-constructor}
\end{equation}
Then every edge has norm $\sqrt{\kappa_d(F)}$,
$H_{\rm geo}=F$, and $S_2^2=16\kappa_d(F)$.
\end{proof}

\begin{remark}[why the determinant is necessary]
It is false in general that
$(\norm{D_1+D_2}+\norm{D_1-D_2})^2$ is at least
$4(\norm{D_1}^2+\norm{D_2}^2)$; nearly parallel diagonals give
counterexamples.  The invariant controlled by the commutator is the Gram
area $q$, and \eqref{eq:gram-chain} is the valid sharp inequality.
\end{remark}

\begin{corollary}[universal three-versus-four switching gain]
\label{cor:gain}
For every nonzero traceless Hermitian target in every finite dimension,
\begin{equation}
 \boxed{\quad
 \frac{\mathcal A_4(F)}{\mathcal A_3(F)}
 =\frac{4}{3\sqrt3}<1.
 \quad}
 \label{eq:gain}
\end{equation}
Equivalently, the fourth kick reduces the optimum squared perimeter by
$1-4/(3\sqrt3)\approx23.02\%$.
On the one-spike cone this becomes
\[
 \mathcal A_3(F)=12\sqrt3\sum_{j=1}^n j b_j,
 \qquad
 \mathcal A_4(F)=16\sum_{j=1}^n j b_j.
\]
The displayed optimal square uses four equal-norm edges and an internal factor
of rank $n$.
\end{corollary}

\section{Relation to earlier records and development provenance}

The dimension-three, dimension-four, and dimension-five inverse-cost papers
already contain the corresponding fixed-dimensional formulas and triangular
loop identities \cite{erikssonQutrit,erikssonFour,erikssonFive}.  The
operational-curvature paper already uses balanced three- and four-kick
commutator loops \cite{erikssonOperational}, and the matrix-isoperimetry paper
establishes the exact four-edge value on the sign-paired spectral face
\cite{erikssonIsoperimetry}.  Accordingly, Theorem~\ref{thm:three} is an
all-dimensional consolidation, while the new scope of Theorem~\ref{thm:four}
is its extension from that equality face to arbitrary traceless Hermitian
targets.

An earlier local, unpublished development package titled \emph{One-Spike
Inverse Self-Commutators in Every Dimension} contains the one-spike formula,
rigidity, and stability lemmas from which this consolidated paper was
developed.  Its hashes and chronology are preserved in the accompanying
priority audit.  A prepared submission sheet is not treated as evidence of
public disclosure.  Public searches of the author's ARR archive,
repositories, and release assets through 30 August 2026 did not locate that exact
manuscript; this negative evidence is not an exhaustive priority guarantee.

The complementary paper \emph{Sharp Rank-Adaptive Bounds for Inverse
Self-Commutators} is now public as ARR-2026-1D2QV1RP1292JREW
\cite{erikssonRankAdaptive}.  It proves the general rank-adaptive bounds, the
uniform one-spike endpoint, both equality classifications, and the exact
sign-cut leakage identity.  It does not contain the arbitrary one-spike
formula, every-optimizer singular rigidity, the sharp stability theorem, or
the all-target kick laws proved here.  The present paper should therefore be
read as an exact interior-cone and synthesis companion, not as a replacement
for that general extremal theorem.

\section{Examples and boundary of the formula}

\paragraph{Qubit.}
For spectrum $(P,-P)$, $\kappa_2(F)=P$ and the formulas give
$\mathcal A_3=12\sqrt3P$ and $\mathcal A_4=16P$.

\paragraph{A zero-padded repeated spectrum.}
For nonzero spectrum $(6,-3,-2,-1)$ in any $d\ge4$,
\[
 \kappa_d=3+2\cdot2+3\cdot1=10,
 \qquad \spec_{>0}(CC^*)=(12,6,2).
\]
Thus every balanced optimum has rank three, while the exact loop costs are
$120\sqrt3$ and $160$.

\paragraph{Why one-spike is a real restriction.}
For $F=\diag(5,1,-3,-3)$, a direct four-level Horn certificate gives
$\kappa_4(F)\ge\lambda_1-\lambda_3=8$.  Equality is attained by
\[
 C=E_{12}+3E_{14}+\sqrt3E_{23}+\sqrt3E_{43},
 \qquad CC^*-C^*C=2F,
 \qquad \frac12\norm C_{\HS}^2=8.
\]
Collapsing the two positive levels into a fictitious spike would instead
produce $3+2\cdot3=9$.  Hence \eqref{eq:one-spike-cost} is not an unrestricted
spectral formula.  The universal loop laws remain valid because they retain
$\kappa_d(F)$ rather than replacing it by the one-spike expression.

\section{Reproducibility, limitations, and conclusion}

The reproduction package contains two independent exact-arithmetic routes.
One checks tail identities, stability bounds and equality families over
rational compositions.  The other checks the weighted-shift and matrix-loop
identities directly, includes nearly parallel diagonals that falsify the
discarded inequality, and verifies the corrected Gram chain.  These programs
test algebraic consequences and fixtures; they do not certify Horn's theorem
or replace the displayed arbitrary-dimensional proofs.

The result is finite-dimensional and uses distinct roles for three notions:
factor product $\kappa_d$, loop perimeter $S_2$, and squared perimeter
$\mathcal A_m$.  No claim is made about quadratic edge energy, operator-norm
action, open paths, coincident-time limits, or higher Magnus orders.  The
priority audit found classical mechanisms close to both halves of the paper
and exact self-overlap in low-dimensional slices.  Novelty therefore remains
a scoped literature-search conclusion rather than a mathematical theorem.

The main structural point is that a difficult inverse matrix invariant can be
read operationally.  On an entire high-rank spectral cone it is an explicit
ordered moment of the negative eigenvalues and rigidifies all optimal
singular data.  Without any spectral restriction, it is exactly the resource
converted by the optimal triangle and square.  The resulting fourth-switch
gain is universal, and the inverse cost is strictly Schur-concave on the
one-spike simplex, even though it is not known in closed form on general sign
patterns.

\section*{Author and AI disclosure}

Lluis Eriksson directed the research, selected the claims, and accepts
responsibility for the manuscript.  OpenAI Codex assisted with algebraic
exploration, literature triage, adversarial proof checking, exact replay,
typesetting, and drafting.  No external peer review or formal proof
certification is claimed.

\small
\bibliographystyle{plain}
\bibliography{references}

\end{document}
